\PassOptionsToPackage{dvipsnames,svgnames,x11names}{xcolor}
\documentclass[tikz,border=8pt]{standalone}
\usepackage{amsmath}
\usepackage{amssymb}
\usepackage{fontawesome5}
\usetikzlibrary{arrows.meta,positioning,calc,shapes,shadows,decorations.pathmorphing,shapes.geometric,backgrounds,decorations.markings,decorations.pathreplacing,decorations.text,fit,shapes.misc,shapes.arrows,matrix,bending,3d,chains,intersections,shapes.symbols,svg.path,shapes.multipart,snakes,shapes.callouts,angles,quotes,fadings,shadows.blur,patterns,patterns.meta}

\providecolor{gold}{RGB}{212,175,55}
\providecolor{silver}{RGB}{192,192,192}
\providecolor{bronze}{RGB}{205,127,50}
\providecolor{darkblue}{RGB}{0,0,139}
\providecolor{lightblue}{RGB}{173,216,230}
\providecolor{darkgreen}{RGB}{0,100,0}
\providecolor{lightgreen}{RGB}{144,238,144}
\providecolor{darkred}{RGB}{139,0,0}
\providecolor{navy}{RGB}{0,0,128}
\providecolor{skyblue}{RGB}{135,206,235}
\providecolor{beige}{RGB}{245,245,220}
\providecolor{cream}{RGB}{255,253,208}

% --- Custom Colors ---
\definecolor{textdark}{RGB}{40, 40, 40}
\definecolor{htmlblue}{RGB}{0, 112, 201}
\definecolor{legacyred}{RGB}{183, 28, 28}
\definecolor{moderngreen}{RGB}{27, 94, 32}

% Code Editor Chrome & Highlighting Colors
\definecolor{codebg}{RGB}{30, 30, 30}
\definecolor{titlebar}{RGB}{55, 55, 55}
\definecolor{linenum}{RGB}{133, 133, 133}
\definecolor{macred}{RGB}{255, 95, 86}
\definecolor{macyellow}{RGB}{255, 189, 46}
\definecolor{macgreen}{RGB}{39, 201, 63}

\definecolor{codeblue}{RGB}{97, 175, 239}
\definecolor{codepurple}{RGB}{198, 120, 221}
\definecolor{codeyellow}{RGB}{229, 192, 123}
\definecolor{codered}{RGB}{224, 108, 117}

\begin{document}
\begin{tikzpicture}[
    show background rectangle,
    background rectangle/.style={fill=white},
    % --- Styles ---
    card/.style={
        rounded corners=10pt, line width=1pt,
        blur shadow={shadow opacity=15, shadow yshift=-3pt, shadow blur radius=6pt},
        inner sep=12pt, text width=6.5cm, align=left, font=\small\sffamily\color{textdark}
    },
    legacycard/.style={
        card, draw=legacyred, top color=legacyred!4, bottom color=white
    },
    moderncard/.style={
        card, draw=moderngreen, top color=moderngreen!4, bottom color=white
    },
    anchorcard/.style={
        card, draw=htmlblue, top color=htmlblue!5, bottom color=white,
        text width=28.36cm, align=center, inner sep=14pt
    },
    arrow/.style={-Stealth, line width=2pt, draw=black!40, rounded corners=6pt},
    greenarrow/.style={-Stealth, line width=2pt, draw=moderngreen, rounded corners=6pt},
    redarrow/.style={-Stealth, line width=2pt, draw=legacyred, rounded corners=6pt}
]

% -------------------------------------------------------------------------
% 1. Titles
% -------------------------------------------------------------------------
\node[font=\Huge\bfseries\sffamily\color{textdark}] at (0, 17.0) {Structural Preservation \& Graceful Degradation};
\node[font=\Large\sffamily\color{gray!80!black}] at (0, 16.0) {Why the \texttt{<picture>} element requires NO JavaScript polyfills};

% -------------------------------------------------------------------------
% 2. Central Code Editor Window (Center X=0, Bounds: (-5.5, 8.5) to (5.5, 13.0))
% -------------------------------------------------------------------------
% Window Frame & Shadow
\fill[codebg, rounded corners=10pt, blur shadow={shadow opacity=25, shadow blur radius=8pt, shadow yshift=-4pt}]
    (-5.5, 8.5) rectangle (5.5, 13.0);
\draw[black!75, line width=1pt, rounded corners=10pt]
    (-5.5, 8.5) rectangle (5.5, 13.0);

% Title Bar
\begin{scope}
    \clip[rounded corners=10pt] (-5.5, 8.5) rectangle (5.5, 13.0);
    \fill[titlebar] (-5.5, 12.4) rectangle (5.5, 13.0);
    \draw[black!60, line width=0.5pt] (-5.5, 12.4) -- (5.5, 12.4);
\end{scope}

% Window Traffic Light Controls
\fill[macred] (-5.0, 12.7) circle (3pt);
\fill[macyellow] (-4.7, 12.7) circle (3pt);
\fill[macgreen] (-4.4, 12.7) circle (3pt);
\node[font=\ttfamily\scriptsize\color{white!60}, anchor=center] at (0, 12.7) {index.html};

% Line 4 Glow / Highlight Box (The Absolute Structural Anchor)
\fill[codeblue!20!codebg, draw=codeblue!45, line width=0.8pt, rounded corners=4pt]
    (-5.3, 9.9) rectangle (5.3, 10.5);

% Code Lines with Line Numbers
% Line 1 (Y = 12.0)
\node[font=\ttfamily\small\color{linenum}, anchor=east] at (-4.7, 12.0) {1};
\node[font=\ttfamily\small, anchor=west] at (-4.4, 12.0) {\textcolor{codeblue}{<picture>}};

% Line 2 (Y = 11.4)
\node[font=\ttfamily\small\color{linenum}, anchor=east] at (-4.7, 11.4) {2};
\node[font=\ttfamily\small, anchor=west] at (-4.4, 11.4) {\hspace*{2.5ex}\textcolor{codeblue}{<source} \textcolor{codepurple}{type}=\textcolor{codeyellow}{"image/avif"} \textcolor{codepurple}{srcset}=\textcolor{codeyellow}{"hero.avif"}\textcolor{codeblue}{>}};

% Line 3 (Y = 10.8)
\node[font=\ttfamily\small\color{linenum}, anchor=east] at (-4.7, 10.8) {3};
\node[font=\ttfamily\small, anchor=west] at (-4.4, 10.8) {\hspace*{2.5ex}\textcolor{codeblue}{<source} \textcolor{codepurple}{media}=\textcolor{codeyellow}{"(min-width: 800px)"} \textcolor{codepurple}{srcset}=\textcolor{codeyellow}{"large.jpg"}\textcolor{codeblue}{>}};

% Line 4 (Y = 10.2) - The Anchor
\node[font=\ttfamily\small\bfseries\color{linenum!70!white}, anchor=east] at (-4.7, 10.2) {4};
\node[font=\ttfamily\small, anchor=west] at (-4.4, 10.2) {\hspace*{2.5ex}\textcolor{codeblue}{<img} \textcolor{codepurple}{src}=\textcolor{codeyellow}{"fallback.jpg"} \textcolor{codepurple}{alt}=\textcolor{codeyellow}{"Hero"}\textcolor{codeblue}{>}};

% Line 5 (Y = 9.6)
\node[font=\ttfamily\small\color{linenum}, anchor=east] at (-4.7, 9.6) {5};
\node[font=\ttfamily\small, anchor=west] at (-4.4, 9.6) {\textcolor{codeblue}{</picture>}};

% -------------------------------------------------------------------------
% 3. Legacy Browser Column (Left, X=-11.0)
% -------------------------------------------------------------------------
\node[font=\Large\bfseries\sffamily\color{legacyred}] at (-11.0, 14.8) {\faInternetExplorer\ Legacy Parsers (e.g., IE 11)};

\node[legacycard] (legacy1) at (-11.0, 12.5) {
    \textbf{\color{legacyred}\faTimesCircle\ Unknown Elements Ignored}\\[4pt]
    Legacy parsers do not recognize \texttt{<picture>} or \texttt{<source>}. Per HTML5 spec, they treat unknown tags as generic inline elements and safely \textbf{parse their children}. No JS errors are thrown; the structural hierarchy is maintained.
};

\node[legacycard] (legacy2) at (-11.0, 9.0) {
    \textbf{\color{legacyred}\faCheckCircle\ Standard Baseline Reached}\\[4pt]
    The parser continues down and natively hits the \texttt{<img>} tag. Recognizing this standard element, it fires a request for \texttt{fallback.jpg} and renders it perfectly.
};

% -------------------------------------------------------------------------
% 4. Modern Browser Column (Right, X=11.0)
% -------------------------------------------------------------------------
\node[font=\Large\bfseries\sffamily\color{moderngreen}] at (11.0, 14.8) {\faChrome\ Modern Parsers (Chrome, Safari)};

\node[moderncard] (modern1) at (11.0, 12.5) {
    \textbf{\color{moderngreen}\faCogs\ Context Wrapper Initialized}\\[4pt]
    Modern parsers recognize \texttt{<picture>} not as a visual element, but as a \textbf{wrapper}. It evaluates the \texttt{<source>} tags for media queries and format support, picking the best match.
};

\node[moderncard] (modern2) at (11.0, 9.0) {
    \textbf{\color{moderngreen}\faCodeBranch\ Property Mutation (\texttt{currentSrc})}\\[4pt]
    The \texttt{<picture>} element \textbf{renders nothing}. Instead, it passes the URL of the winning \texttt{<source>} directly to the \texttt{<img>} tag, mutating its \texttt{currentSrc} property.
};

% -------------------------------------------------------------------------
% 5. The Bottom "Anchor" Banner (Center X=0, Y=4.5, W=24.0cm)
% -------------------------------------------------------------------------
\node[anchorcard] (anchor) at (0, 4.5) {
    \Large\textbf{\color{htmlblue}\faAnchor\ The \texttt{<img>} Tag is the Absolute Structural Anchor}\\[6pt]
    \normalsize Without the \texttt{<img>} tag, a \texttt{<picture>} element is entirely invisible in \textbf{both} modern and legacy browsers. \\ The \texttt{<picture>} tag provides the cascading logic, but the \texttt{<img>} tag paints the actual pixels to the DOM.
};

% -------------------------------------------------------------------------
% 6. Precision Orthogonal Routing & Grouping (Braces & Connectors)
% -------------------------------------------------------------------------

% Legacy (Left) Routing:
% Grouping brace for lines 1-3 spanning (-5.9, 10.6) to (-5.9, 12.2)
\draw[decorate, decoration={brace, amplitude=6pt}, legacyred, line width=1.5pt] (-5.9, 10.6) -- (-5.9, 12.2);
% Upper dashed arrow from brace center (-6.15, 11.4) to legacy1.east
\draw[redarrow, dashed] (-6.15, 11.4) -- (-6.8, 11.4) |- (legacy1.east);
% Lower solid arrow from line 4 coordinate (-5.9, 10.2) to legacy2.east
\draw[redarrow] (-5.9, 10.2) -- (-6.8, 10.2) |- (legacy2.east);

% Modern (Right) Routing:
% Grouping brace for lines 1-3 spanning (5.9, 12.2) to (5.9, 10.6)
\draw[decorate, decoration={brace, amplitude=6pt}, moderngreen, line width=1.5pt] (5.9, 12.2) -- (5.9, 10.6);
% Upper dashed arrow from brace center (6.15, 11.4) to modern1.west
\draw[greenarrow, dashed] (6.15, 11.4) -- (6.8, 11.4) |- (modern1.west);
% Lower solid arrow from line 4 coordinate (5.9, 10.2) to modern2.west
\draw[greenarrow] (5.9, 10.2) -- (6.8, 10.2) |- (modern2.west);

% Downward Resolution Arrows to Anchor Banner
\draw[redarrow] (legacy2.south) -- (legacy2.south |- anchor.north);
\draw[greenarrow] (modern2.south) -- (modern2.south |- anchor.north);
\draw[arrow, dashed] (0, 8.1) -- (anchor.north);

\end{tikzpicture}
\end{document}
